% \chapter{结\texorpdfstring{\quad}{}论}
\chapter{总结与展望}

涵道风扇式无人机因空气动力效率高、安全性高和可垂直起降等优势，得到了国内外相关研究机构的重视。但涵道的存在引起了气流与涵道之间的强相互耦合作用，对DFUAV的自主飞行控制提出了更高的要求。因此，本文针对DFUAV的飞行控制和轨迹规划问题展开研究。主要的工作及结论如下：

(1)基于牛顿-欧拉方法建立了DFUAV的飞行控制刚体模型。模型中使用地面坐标系和机体坐标系的多坐标系方法表示DFUAV的平移运动与旋转运动关系，并从涵道风扇、机身、涵道环翼、控制舵面和固定气动面等多个方面分析了作用在机体上的力与力矩，定量地分析了各类作用力与力矩对DFUAV的动态特性的影响。飞行控制刚体模型的建立为后续控制器的设计提供了理论基础。

(2)自主搭建了一架试验样机用于实际飞行实验。机载航电系统搭载的微控制器和各种传感器器件能够满足实验中需要的的算力和数据质量等要求。此外，在MATLAB环境下建立了DFUAV的仿真模型，该模型基于龙格-库塔法的ode45数值求解器求解DFUAV的非线性动态方程，相同输入下，仿真结果与实际飞行结果吻合较好。

(3)针对DFUAV的姿态控制问题，提出了基于INDI与优先级控制分配的理论框架与工程实现方案。首先对姿态动力学进行合理简化，将外部力矩分为可控力矩和不可控力矩。在INDI姿态控制架构下，根据期望姿态计算出反馈输入力矩，根据陀螺力矩机理使用控制舵面偏转进行补偿，并且通过机载陀螺仪对不可控力矩在每个周期内实时补偿得到INDI输入力矩。结合两部分输入，采用优先级控制分配算法对输入进行合理分配，以最大范围返回容许控制。最后通过一系列仿真实验与实际飞行实验证明了INDI姿态控制方法相比PID方法在姿态跟踪性能和抗干扰性能方面表现更好，并且从分配误差的角度证明了优先级控制分配算法相比伪逆方法误差更小。

(4)针对DFUAV的位置控制问题，在姿态控制的基础上提出了基于INDI的位置控制架构，形成外环-内环协同的扰动观测补偿机制。首先对位置动力学模型作一定的简化，统一为加速度量纲，便于基于加速度计测量值估计机体受力。在INDI位置控制架构下，根据期望位置指令以及加速度计测量值计算出所需推力增量，每个周期内的推力增量中已经包含了对当前周期内机体所受扰动的补偿。得到推力在各轴的分量之后，采用几何解算方法计算出期望的姿态指令输入内环控制器。最后通过一系列仿真实验与实际飞行实验证明了INDI位置控制方法相比PID方法在抗干扰性能方面表现出色，并且INDI位置控制器在轨迹跟踪能力上同样表现较好。

(5)针对DFUAV的轨迹规划问题，提出了结合C-RRT*算法和B样条曲线的轨迹规划方法。首先介绍了RRT算法、RRT-Connect算法和RRT*算法的基本原理，并分析了各自的优势与不足之处，然后通过结合RRT-Connect算法和RRT*算法的优势提出了C-RRT*算法，并且采用冗余节点进一步优化路径，减小弯曲度。针对C-RRT*算法搜索的初始路径不平滑的问题，利用B样条曲线的性质对初始路径平滑处理，并对路径赋予满足动力学约束的时间参数。最后通过仿真实验对比验证了C-RRT*算法综合性能比RRT-Connect和RRT*算法更优，并且B样条曲线也能够平滑路径并得到满足动力学约束的速度和加速度曲线，轨迹跟踪结果进一步表明了该方法的有效性。

本文针对DFUAV提出的飞行控制方案与轨迹规划方案取得了一定的成果，可供同行作参考借鉴。但是也清晰地认识到本文的工作仍有不足之处，据此对未来的研究工作作出展望：
% 为便于实现基于INDI的姿态控制器设计，方案中对姿态动力学模型做了一定程度的简化，但也因此丢失了模型的部分细节信息，影响了DFUAV的动态性能。此外，在保留计算效率的同时也能保证系统的动态响应能力。

(1)DFUAV的部分系统辨识参数以及控制参数在系统各模态下可能有不同的取值，采用单一的参数应对系统的不同模态会限制系统性能的发挥。未来可以考虑采用自适应方案，针对不同环境下无人机的状态来动态调整辨识的参数和控制器参数，以进一步提升系统的鲁棒性和动态响应能力。

(2)受算力和存储空间等限制，现有的飞行平台还不足以在线规划出从起点到目标点的无碰撞平滑路径，所以本文的轨迹规划与跟踪实验仅在仿真环境下开展。未来将考虑采用具有更高算力与存储空间的飞行平台做在线规划。

(3)本文的轨迹规划方法仅考虑了二维平面内障碍物静止的轨迹规划问题，未来将考虑处于三维空间且存在动态障碍物的场景，以扩大DFUAV的应用范围。

